\originalsection{体源与面源从何而来}
设极化介质占据区域 $\Omega$, $\vct P$ 在内部光滑. 偶极势相加给
$V_b=(4\pi\eps)^{-1}\int_\Omega\vct P(\rr')\cdot\nabla'[1/|\rr-\rr'|]\dd^3x'$.
注意对源坐标求梯度的符号为正, 对观察坐标求梯度则为负. 对 $\nabla'\cdot(\vct P/R)$ 积分,
\[
 V_b=\frac1{4\pi\eps}\left[
 \oint_{\partial\Omega}\frac{\vct P\cdot\nn}{R}\dd a'
 -\int_\Omega\frac{\nabla'\cdot\vct P}{R}\dd^3x'\right].
\]
与普通电荷的库仑势逐项比较, 才能认出 $\rho_b=-\nabla\cdot\vct P$, $\sigma_b=\vct P\cdot\nn$. 把体源积分和面源相加, 总束缚电荷为零, 因为它们来自原来电中性的偶极而非添入净电荷. 若介质只是更大极化系统的一部分, 则须把人工切面的另一侧也计入, 不能把切面电荷误当整系统的新电荷.

以 $\vct P=P_0\hat{\vct z}$ 的均匀极化球为例, 体源为零, 球面 $\sigma_b=P_0\cos\theta$. 表面仅有 $\ell=1$ 模态, 令内势 $Ar\cos\theta$, 外势 $Ba^3\cos\theta/r^2$. 连续势给 $A=B$, 法向场差为 $(2B-(-A))\cos\theta=3A\cos\theta$, 所以 $A=P_0/(3\eps)$. 球内电场 $-\vct P/(3\eps)$, 球外等效偶极矩为体积乘 $\vct P$. 这个负的退极化内场来自两极的束缚电荷, 不表示材料极化率本身为负.

$\vct P$ 是响应, $\vct D$ 是把自由源留在高斯定律中的辅助场. 即便 $\rho_f=0$, 也只能推出 $\nabla\cdot\vct D=0$, 不能推出 $\vct D=0$ 或 $\ee=0$. 比如均匀外场中的介质球, 远处非零场是边界数据, 第5章例题正是无自由体电荷却有非零内部场的情况. 在没有球、柱或平面等足够对称性的几何中, 不能仅从一个高斯曲面的通量决定局部 $\vct D$.
